% ID: SEL-D0059-OA-E01
% ESTADO: cotejada. Segunda pasada visual y borrador compilado revisados.
% Fuente: Selectividad/sel_2011_matematicas_aplicadas/222-EXAMEN 2.pdf
% SHA-256: 998f20402054ca6d0454da83ffc8e1ff6951a576b48836a32e81e524c5985214
% Página: 1; opción A; ejercicio original 1.

\begin{enumerate}[label=\alph*)]
\item \textbf{(1.25 puntos)} Dada la matriz $A=\begin{pmatrix}1&5&6\\0&1&7\\0&0&1\end{pmatrix}$, calcule $(I_3-A)^3$.
\item \textbf{(1.25 puntos)} Dadas las matrices $B=\begin{pmatrix}1&a\\b&3\end{pmatrix}$, $C=\begin{pmatrix}-1\\3\end{pmatrix}$, $D=\begin{pmatrix}5\\10\end{pmatrix}$, determine $a$ y $b$ de manera que $B\cdot C-D=O$, siendo $O$ la matriz nula.
\end{enumerate}
